% macros.tex – Einheitliche Präambel für QR-Theorie v6.00
% Optimiert für pdfLaTeX & GitHub Actions
% Keine \documentclass hier; steht in main.tex.

% ---------- Encoding & Sprache ----------
\usepackage[utf8]{inputenc}
\usepackage[T1]{fontenc}
\usepackage{lmodern}
\usepackage[ngerman,english]{babel}
\usepackage{microtype}
\usepackage{csquotes}

% ---------- Seitenlayout ----------
\usepackage[margin=2.5cm,bindingoffset=0.5cm]{geometry}
\usepackage{fancyhdr}
\pagestyle{fancy}
\fancyhf{}
\fancyhead[LE,RO]{\thepage}
\fancyhead[LO]{\nouppercase{\rightmark}}
\fancyhead[RE]{\nouppercase{\leftmark}}
\setlength{\headheight}{14.5pt}
\addtolength{\topmargin}{-2.5pt}

% ---------- Mathematik ----------
\usepackage{amsmath,amssymb,amsthm,mathtools}
\usepackage{physics}
\numberwithin{equation}{section}

% ---------- Abbildungen & Tabellen ----------
\usepackage{graphicx}
\usepackage{booktabs}
\usepackage{array}
\usepackage{multirow}
\usepackage{float}
\usepackage{caption}
\usepackage{subcaption}

% ---------- Einheiten & Zahlen (siunitx v3) ----------
\usepackage{siunitx}
\sisetup{
  uncertainty-mode = separate,
  per-mode = symbol,
  range-phrase = {---},
  list-final-separator = {, and }
}
\AtBeginDocument{\RenewCommandCopy\qty\SI} % nur 1x!
\DeclareSIUnit\parsec{pc}
\DeclareSIUnit\astronomicalunit{au}
\DeclareSIUnit\lightyear{ly}
\DeclareSIUnit\Mpc{Mpc}
\DeclareSIUnit\Gpc{Gpc}

% ---------- Farben ----------
\usepackage{xcolor}
\definecolor{qrblue}{RGB}{0,102,204}
\definecolor{qrred}{RGB}{204,0,51}
\definecolor{qrgreen}{RGB}{0,153,76}

% ---------- Hyperlinks & Referenzen ----------
\usepackage[colorlinks=true,linkcolor=qrblue,citecolor=qrred,urlcolor=qrgreen]{hyperref}
\usepackage{cleveref}

% ---------- Bibliographie ----------
\usepackage[backend=biber,style=authoryear-comp,doi=false,isbn=false,url=true,maxbibnames=3]{biblatex}
\addbibresource{bibliography/qr_references.bib}

% ---------- Algorithmen & Code ----------
\usepackage{algorithm}
\usepackage{algorithmic}
\usepackage{listings}
\lstset{
  basicstyle=\footnotesize\ttfamily,
  breaklines=true,
  frame=single,
  numbers=left,
  numberstyle=\tiny,
  captionpos=b
}

% ---------- Spezielle Physik-Pakete ----------
\usepackage{tensor}
\usepackage{slashed}

% ---------- Grafikpfade ----------
\graphicspath{{figures/}{./}}

% ---------- QR-Theorie spezifische Makros ----------
\newcommand{\QR}{\text{QR}}
\newcommand{\LCDM}{\Lambda\text{CDM}}
\newcommand{\Inu}{I_\nu}
\newcommand{\lpoi}{\ell_\pi}
\newcommand{\mathcalI}{\mathcal{I}}
\newcommand{\mathcalO}{\mathcal{O}}
\newcommand{\mathcalW}{\mathcal{W}}
\newcommand{\mathcalS}{\mathcal{S}}
\newcommand{\mathcalF}{\mathcal{F}}

% Mathematische Operatoren
\DeclareMathOperator{\diag}{diag}
\DeclareMathOperator{\sign}{sign}

% Theorem-Umgebungen
\theoremstyle{definition}
\newtheorem{axiom}{Axiom}[section]
\newtheorem{definition}{Definition}[section]
\newtheorem{theorem}{Theorem}[section]
\newtheorem{proposition}{Proposition}[section]
\newtheorem{corollary}{Corollary}[theorem]
\newtheorem{lemma}[theorem]{Lemma}

\theoremstyle{remark}
\newtheorem{remark}{Remark}[section]
\newtheorem{example}{Example}[section]

% QR-Notation
\newcommand{\qroperator}[1]{\mathcal{O}^{(#1)}}
\newcommand{\projweight}[1]{\mathcal{W}(#1)}
\newcommand{\infodensity}[1]{I_\nu(#1)}
\newcommand{\loginfo}[1]{\varphi(#1)} % konsistent zur goldenen Ratio
\newcommand{\goldenratio}{\varphi}
\newcommand{\scalehierarchy}[1]{\lambda_{#1}}

% Einheiten-Makros
\newcommand{\kmsmpc}{\si{\kilo\metre\per\second\per\Mpc}}
\newcommand{\arcmin}{\si{\arcminute}}

% QR-Gleichungen (Symbolik)
\newcommand{\qrfieldequation}{G_{\mu\nu} + \Lambda g_{\mu\nu} = \frac{8\pi G}{c^4} T_{\mu\nu}^{\text{eff}}}
\newcommand{\qrstressenergy}{T_{\mu\nu}^{\QR} = \frac{c^4}{8\pi G} \left[ (\log \Inu)^2 \varphi^2 g_{\mu\nu} + \nabla_\mu \varphi \nabla_\nu \varphi - \frac{1}{2} g_{\mu\nu} \nabla_\alpha \varphi \nabla^\alpha \varphi \right]}

% ---------- Titelseiten ----------
\newcommand{\qrtitle}{The QR Theory v6.00: A Complete Synthesis of Quantum Mechanics and General Relativity through Projective Information Dynamics}
\newcommand{\qrauthor}{Frank Kannstädter}
\newcommand{\qraffiliation}{Independent Researcher in Theoretical Physics, Frankfurt am Main, Germany}

% Abstract & Keywords
\newenvironment{qrabstract}
{\begin{quotation}\noindent\textbf{Abstract:}\space}
{\end{quotation}}
\newcommand{\qrkeywords}[1]{\noindent\textbf{Keywords:} #1}

% Inhaltsverzeichnis-Tiefe
\setcounter{tocdepth}{3}
\setcounter{secnumdepth}{3}

% Beschriftungen
\captionsetup{
  font=small,
  labelfont=bf,
  format=plain,
  justification=centering
}

% QR-Axiome Umgebung
\newenvironment{qraxiom}[1]
{\begin{axiom}[#1]\itshape}
{\end{axiom}}

% Validierungstabellen
\newcommand{\validationtable}[4]{%
\begin{table}[H]
\centering
\caption{#1}
\label{#2}
\begin{tabular}{#3}
#4
\end{tabular}
\end{table}
}

% Residuenplots
\newcommand{\residualplot}[2]{%
\begin{figure}[H]
\centering
\includegraphics[width=0.8\textwidth]{#1}
\caption{#2}
\label{fig:residual}
\end{figure}
}

% Warnungen filtern
\usepackage{silence}
\WarningFilter{siunitx}{Option "separate-uncertainty"}
\WarningFilter{latex}{You have requested package}

% TODO-Notizen
\usepackage{todonotes}
\newcommand{\qrtodo}[1]{\todo[color=yellow!40]{QR: #1}}

% ---------- Finale Einstellungen ----------
\raggedbottom
\frenchspacing